\section{向量空间的对偶}\label{sec:vector-space-duality}

本节介绍从内在角度研究切空间所需的向量空间对偶的关键概念. 讨论最后将得到一个性质: 每个有限维向量空间都与它的双重对偶典范同构.

先定义向量空间的对偶. 它由该向量空间到基域的所有线性映射组成.

\begin{Definition}{向量空间的对偶}\leavevmode\label{def:dual-vector-space}
设 $V$ 为向量空间. $V$ 的\emph{对偶}记为 $V^{\vee}$, 是所有线性函数 $\phi:V\to K$ 的集合.
\end{Definition}

注意, 在到基域 $K$ 的函数通常的加法和标量乘法下, $V^{\vee}$ 本身也是向量空间 (Exercise~\ref{ex:dual-is-vector-space}). 此外, 给定 $V$ 的一组基, $V^{\vee}$ 上有一组自然的\textquotedbl{}对偶\textquotedbl{}基. 下面的定义基于这一事实: 任意线性函数 $V\to K$ 由它在 $V$ 的一组基的各元素上的取值唯一确定.

\begin{Definition}{一组基的对偶}\leavevmode\label{def:dual-of-a-basis}
若 $V$ 为有限维向量空间, 基 $\{v_{1},\ldots,v_{n}\}$ 的大小为 $n$, 则 $\{v_{1},\ldots,v_{n}\}$ 的\emph{对偶}是函数集 $\{v_{1}^{\vee},\ldots,v_{n}^{\vee}\}\subseteq V^{\vee}$, 它由下式确定:
\[
  v_{i}^{\vee}(v_{j})=
  \begin{cases}
    1 & i=j,\\
    0 & i\neq j.
  \end{cases}
\]
\end{Definition}

来看一个具体例子.

\begin{Example}{$K^{2}$ 的一组基的对偶}\leavevmode\label{ex:dual-of-a-basis-of-k-squared}
设 $V=K^{2}$, 考虑由以下两个向量组成的基:
\[
  v_{1}=(1,1),\qquad v_{2}=(1,-1).
\]
注意 $v_{1}^{\vee}:V\to K$ 将 $(1,1)$ 映到 $1$, 将 $(1,-1)$ 映到 $0$, 而 $v_{2}^{\vee}$ 将 $(1,1)$ 映到 $0$, 将 $(1,-1)$ 映到 $1$. 因此, 给定任意 $(a,b)\in K^{2}$, 写成
\[
  (a,b)=\frac{a+b}{2}(1,1)+\frac{a-b}{2}(1,-1)
\]
便得到
\[
  v_{1}^{\vee}(a,b)=\frac{a+b}{2}
  \qquad\text{和}\qquad
  v_{2}^{\vee}(a,b)=\frac{a-b}{2}.
\]
\end{Example}

正如可以预期的那样, $V$ 的一组基的对偶是 $V^{\vee}$ 的一组基. 现在验证这一点.

\begin{Proposition}{一组基的对偶是一组基}\leavevmode\label{prop:dual-basis-is-a-basis}
若向量空间 $V$ 有一组大小为 $n$ 的基 $\{v_{1},\ldots,v_{n}\}$, 则 $\{v_{1}^{\vee},\ldots,v_{n}^{\vee}\}$ 是 $V^{\vee}$ 的一组基.
\end{Proposition}
\begin{Proof}
证明 $\{v_{1}^{\vee},\ldots,v_{n}^{\vee}\}$ 线性独立且张成 $V^{\vee}$.

为证明线性独立, 设对某些 $a_{1},\ldots,a_{n}\in K$, 有 $a_{1}v_{1}^{\vee}+\cdots+a_{n}v_{n}^{\vee}=0$. 在 $v_{i}$ 处取值, 得
\[
  0=(a_{1}v_{1}^{\vee}+\cdots+a_{n}v_{n}^{\vee})(v_{i})=a_{i}.
\]
因此 $a_{1}=\cdots=a_{n}=0$, 这说明 $\{v_{1}^{\vee},\ldots,v_{n}^{\vee}\}$ 线性独立.

为证明 $\{v_{1}^{\vee},\ldots,v_{n}^{\vee}\}$ 张成 $V^{\vee}$, 设 $\phi\in V^{\vee}$, 即 $\phi:V\to K$ 为线性函数. 考虑线性组合 $\phi(v_{1})v_{1}^{\vee}+\cdots+\phi(v_{n})v_{n}^{\vee}$. 注意对每个 $j=1,\ldots,n$, 由 $v_{1}^{\vee},\ldots,v_{n}^{\vee}$ 的定义有
\[
  (\phi(v_{1})v_{1}^{\vee}+\cdots+\phi(v_{n})v_{n}^{\vee})(v_{j})=\phi(v_{j}).
\]
由于 $\{v_{1},\ldots,v_{n}\}$ 是 $V$ 的一组基, 且线性函数 $\phi(v_{1})v_{1}^{\vee}+\cdots+\phi(v_{n})v_{n}^{\vee}$ 和 $\phi$ 在每个基元素上相同, 可知 $\phi=\phi(v_{1})v_{1}^{\vee}+\cdots+\phi(v_{n})v_{n}^{\vee}$. 由于 $\phi$ 是任意的, 这说明 $\{v_{1}^{\vee},\ldots,v_{n}^{\vee}\}$ 张成 $V^{\vee}$.
\end{Proof}

\BookPairedNote
  {前一个结论蕴含, 有限维向量空间的对偶本身也是同维的有限维向量空间. 特别地, 给定 $V$ 的一组基, 就有一个同构 $V\cong V^{\vee}$, 它将每个基元素映到相应的对偶基元素.}
  {注意 Proposition~\ref{prop:dual-basis-is-a-basis} 对无限维向量空间不成立, 见 Exercise~\ref{ex:infinite-dimensional-dual-basis-failure}.}
由此产生一个重要问题: 这个同构是否与基无关? 答案是\textquotedbl{}否\textquotedbl{}, 如下例所示.

\begin{Example}{不同的基定义不同的同构 $V\cong V^{\vee}$}\leavevmode\label{ex:duality-isomorphisms-depend-on-basis}
设 $V=K^{2}$, 考虑分别由下面两对元素给出的两组基:
\[
  v_{1}=(1,0),\quad v_{2}=(0,1)
  \qquad\text{和}\qquad
  w_{1}=(1,0),\quad w_{2}=(1,1).
\]
可以验证 $\{v_{1},v_{2}\}$ 的对偶基由
\[
  v_{1}^{\vee}(a,b)=a\qquad\text{和}\qquad v_{2}^{\vee}(a,b)=b
\]
给出, 而 $\{w_{1},w_{2}\}$ 的对偶基由
\[
  w_{1}^{\vee}(a,b)=a-b\qquad\text{和}\qquad w_{2}^{\vee}(a,b)=b
\]
给出. 特别地, 由于 $v_{1}=w_{1}$ 但 $v_{1}^{\vee}\neq w_{1}^{\vee}$, 可知把 $v_{1},v_{2}$ 映到 $v_{1}^{\vee},v_{2}^{\vee}$ 的 $V$ 与 $V^{\vee}$ 之间的同构, 不同于把 $w_{1},w_{2}$ 映到 $w_{1}^{\vee},w_{2}^{\vee}$ 的同构.
\end{Example}

我们已经看到如何选择一组基来构造 $V$ 与 $V^{\vee}$ 之间的同构. 但前例表明, 对不同的基, 这些同构通常不同. 事实上, 无论怎么尝试, 如果不作一些额外选择 (比如选定一组基), 就无法构造非零向量空间与其对偶之间的同构.

\BookPairedNote
  {那么, 以下事实或许有点出人意料: 再取一次对偶, 就会发现 $V$ 与 $(V^{\vee})^{\vee}$ 之间存在典范同构.}
  {通常称 $(V^{\vee})^{\vee}$ 为 $V$ 的\textquotedbl{}双重对偶\textquotedbl{}.}

\begin{Theorem}{$V=(V^{\vee})^{\vee}$}\leavevmode\label{thm:vector-space-canonically-isomorphic-to-double-dual}
若 $V$ 是有限维向量空间, 则 $V$ 与 $(V^{\vee})^{\vee}$ 之间存在典范同构.
\end{Theorem}
\begin{Proof}
定义 $\omega:V\to(V^{\vee})^{\vee}$, 把 $V$ 中向量 $v$ 送到 $V^{\vee}$ 上相应的\textquotedbl{}取值函数\textquotedbl{}, 其定义为
\[
  \begin{aligned}
    e_{v}:V^{\vee}&\to K,\\
    \phi&\mapsto\phi(v).
  \end{aligned}
\]
稍加思考即可看出, 对任意 $a\in K$ 及 $v,w\in V$, 有 $e_{av+w}=ae_{v}+e_{w}$. 换句话说, $\omega$ 是线性映射. 注意, 线性映射 $\omega$ 的确是典范地定义的, 因为其定义不需要任何额外选择 (比如选择一组基).

还须证明 $\omega$ 是同构, 为此选取一组基. 设 $V$ 的维数为 $n$, 且 $\{v_{1},\ldots,v_{n}\}\subseteq V$ 为 $V$ 的一组基. 在对偶基元素 $v_{j}^{\vee}\in V^{\vee}$ 处对取值函数 $e_{v_{i}}$ 取值, 得
\[
  e_{v_{i}}(v_{j}^{\vee})=v_{j}^{\vee}(v_{i})=
  \begin{cases}
    1 & i=j,\\
    0 & i\neq j.
  \end{cases}
\]
换句话说, 这表明 $\{e_{v_{1}},\ldots,e_{v_{n}}\}$ 是对偶基 $\{v_{1}^{\vee},\ldots,v_{n}^{\vee}\}$ 的对偶基, 因而由 Lemma~\ref{prop:dual-basis-is-a-basis} 可知它是一组基. 由于 $\omega$ 是把 $V$ 的一组基送到 $(V^{\vee})^{\vee}$ 的一组基的线性映射, 由标准的线性代数论证得 $\omega$ 是同构.
\end{Proof}

\BookPairedNoteLeft
  {熟悉了向量空间的对偶 (和双重对偶) 后, 下一节回到切空间的话题. 那里将给出切空间的一种内在描述, 它不依赖于仿射簇置于仿射空间中的方式.}
  {注意, 对任意向量空间 $V$, $\omega:V\to(V^{\vee})^{\vee}$ 都是单射, 但只有在 $\dim(V)<\infty$ 时才是满射.}

\BookExerciseHeading
\begin{BookExercises}
  \item\label{ex:dual-is-vector-space} 设 $V$ 为向量空间. 证明 $V^{\vee}$ 也是向量空间.
  \item\label{ex:distinct-bases-may-give-same-duality-isomorphism} 给出 $K^{2}$ 的两组不同的基 $\{v_{1},v_{2}\}$ 和 $\{w_{1},w_{2}\}$, 使得由 $v_{i}\mapsto v_{i}^{\vee}$ 和 $w_{i}\mapsto w_{i}^{\vee}$ 给出的同构 $V\to V^{\vee}$ 相同.
  \item\label{ex:infinite-dimensional-dual-basis-failure} 考虑无限维向量空间 $K[x]$ 及其基 $\{1,x,x^{2},x^{3},\ldots\}$. 证明线性映射
  \[
    \begin{aligned}
      K[x]&\to K,\\
      f&\mapsto f(1)
    \end{aligned}
  \]
  不在 $1^{\vee},x^{\vee},(x^{2})^{\vee},(x^{3})^{\vee},\ldots$ 的张成空间中.
  \item\label{ex:linear-maps-from-base-field} 设 $V$ 为向量空间, $W$ 为线性映射 $K\to V$ 的集合. 证明 $V$ 和 $W$ 是典范同构的向量空间.
  \item\label{ex:linear-maps-into-finite-dimensional-vector-space} 设 $V$ 和 $W$ 为向量空间, $W$ 的有限维数为 $n$, 且 $\{w_{1},\ldots,w_{n}\}$ 为 $W$ 的一组基. 构造 $(V^{\vee})^{n}$ 与线性映射 $V\to W$ 所成向量空间之间的一个同构.
  \item\label{ex:dual-and-double-dual-of-linear-map} 设 $V$ 和 $W$ 为向量空间. 给定线性映射 $\phi:V\to W$, 定义 $\phi$ 的\emph{对偶}为
  \[
    \begin{aligned}
      \phi^{\vee}:W^{\vee}&\to V^{\vee},\\
      \psi&\mapsto\psi\circ\phi.
    \end{aligned}
  \]
  \begin{enumerate}
    \item 说明对任意线性映射 $\phi:V\to W$ 和 $\psi\in W^{\vee}$, 为什么有 $\psi\circ\phi\in V^{\vee}$.
    \item 证明 $\phi^{\vee}$ 是线性映射.
    \item 设 $V$ 和 $W$ 均为有限维. 证明 $(\phi^{\vee})^{\vee}=\phi$.
  \end{enumerate}
\end{BookExercises}
